\documentclass[CSHFoundation.tex]{subfiles}
\begin{document}

\chapter{Foundation}
\centerline{\includegraphics[scale=0.35]{6-Flogo.png}}
\section{Council Constitution}
\subsection{Purpose}
The council exists to develop and continue relationship between the school, church, and home.\footnote{It is a autonomous relationship in itself, and will not be directly governed by any one organization or group}
\subsection{Layout}
It will consist of 12 members\footnote{12 is the number used in the Bible, but it also represents one of the best numbers to serve as a base in math. 12 can be divided by 1,2,3,4,6,and 12. Unlike 10, this provides an motion process that can can require 75\% of the votes}

\begin{itemize}
 \item 4 elected members from the church
 \item 4 elected members from the school
 \item 4 elected members from the homes
 \end{itemize} 

Each member will be elected to serve for 4 years, for a maximum of 12 years.\footnote{2 years is not enough time to properly make long term decisions. A 4 year term gaurentees that long term decisions are made by those with experience.}

\vspace*{4mm}
Additionally a chair will be elected for 2 years, from the existing members of the council\footnote{Members who are in their $3^{rd}$ year may not run}

\vspace*{4mm}
Each year there will be an election for one new member from each group of the council\footnote{This ensures a turn over rate of 3 and no more per year. This also stagers the elections, so that at any given year, there are not all new members belonging to an group}

\subsection{Procedures}

\begin{itemize}
\item All members must belong to a local church
\item 9 votes are required to pass a motion (75\%)
\item Meetings are open for the public to view
\item Meetings will be digitally recorded
\item Leaders who are living in sin must be confronted one on one, and then by the entire council, and lastly in their respective local church
\end{itemize} 



\section{Beliefs}

This documents fundamental beliefs that led the the inclusion of practises. If you have a problem with a practise, the solution is not to first dismiss the practise, but to understand the belief below that practise. If you disagree with the practise, but agree with the belief, then you are in good company. It does not matter if you change the practise, if the belief is upheld.

All of our practises, all of our interactions, can only be properly understood when we understand who God is.

\subsubsection{God}

\begin{enumerate}
\item Omniscient
\item Omnipresent
\item Omnipotent
\item Immortal
\item Relational
\item Father, Son, Holy Spirit
\item Holy
\item Life
\item Love \footnote{These lists and ideas come from an understanding of the God portrayed in both the pre-Christ, Christ, and post-Christ Testament (also known as the Old Testament, the Gospels, and the New Testament}
\end{enumerate}

The interaction between people and knowledge need not be divorced from each other. The way we interact with God is the same way we should interact with each other, and with knowledge. \footnote{Our relationship to God is the most fundamental relationship. It is through that relationship that all other relationships flow}



\subsubsection{Knowledge Sphere}



\begin{enumerate}
\item Knowledge requires exploration
\item Knowledge requires real examples
\item Knowledge requires giving of ourselves
\item Knowledge requires commitment to learn
\item Knowledge requires communication
\item Knowledge requires submission to truth
\item Knowledge requires errors and correction
\item Knowledge requires application to life
\item Knowledge requires love \footnote{In both application, but also in how we approach it. If we do not love knowledge, we will not apply it to our lives}
\end{enumerate}



\subsubsection{Social Sphere}



\begin{enumerate}
\item Relationship requires exploration
\item Relationship requires physical interaction
\item Relationship requires giving of ourselves
\item Relationship requires commitment
\item Relationship requires communication
\item Relationship requires submission to each other
\item Relationship requires change and repentance
\item Relationship requires living life together
\item Relationship requires love
\end{enumerate}

\subsection{Policies}

\subsubsection{Avoid Abortion, Adopt Always}

Any child in need of a home will be adopted, if it will save their life from being aborted. Steps will be taken to ensure that the baby is not killed by the mother



\subsubsection{Zero Prior Knowledge:}

Given a standard set of instruction on basic computer usage, a person should be able to sit down at any new piece of software, and without any help, use it



\subsubsection{Face2Face Policy:}

The most effective way to produce change in another person is through the use of face to face conversation



\subsubsection{Open Hearts Policy:}

Our hearts are open to those who cry out for help. They will not be turned away open-handed. Either with physical, emotional, intellectual, or spiritual help. As far as we have means, we help others. There is no child who is beyond our help. There is no difficulty we are not willing to tackle.


\subsubsection{Profit}
Knowledge must always be given without pay. Relationships must be build with no pretext of monetary gain. To this end, all three structures are strictly non-profit, and will fall under immediate investigation with profits are being made. In the case of the school and church, those who work should receive wages for their work, but this money should not come from students, but rather from the community who donates monthly to the school. Likewise, no profit should ever be taken in order for someone to receive the Gospel (or any biblical instruction). The wages come from those who believe in the cause, and donate their money to finance the people involved.

\subsubsection{Homosexuality}
The CSH Foundation recognizes that those who struggle with homosexuality believe that they will find peace in expressing what they feel inside. The CSH Foundation also recognizes that one's identity and worth are not found in what someone feels about themselves, but in their relationship to God. As such, students who feel attracted to the same sex are encouraged to come to all spheres, but they do so under the understanding that the community is there to help them find their identity in Christ.\footnote{We were created male and female. While Galatians 3:28 speaks of our belonging in Christ, it does not speak to destroying all roles. Paul was not advocating freedom for every slave, but rather to work in their role. Homosexuality is a hard road to follow, but we believe that true peace, and true joy, and true love, are found in the will of God alone. More could be argued against homosexuality from an scientific, sociological, or logical standpoint, but that is left as an exercise for the reader.}

\subsection{Women in Leadership}

Women are not permitted to teach or have authority over men, in any sphere. \footnote{This is not to say that the any sphere can exist without women. It cannot. The home cannot exist with women, schools cannot exist very well without women, and churches desperately need more females to help. This is not meant to destroy the ministry of women, but rather to channel it. Teach means more then just to spout knowledge, but rather to have a relationship where they are actively guiding another individual. Boys transition into men when they come of age to make their own decisions in life. In the Canadian culture, this is mid 20's or so, when a man becomes independent from his parents}

\section{Social}

Each individual is composed of knowledge, built on the core of God (Father, Son, Holy Spirit).

	In the church, we function as a community founded on the knowledge of God. As God is our Creator, and the foundation of all knowledge, the church unit is the more fundamental organization in society. When the Church fails, society as a whole fails.
	
	Families form the emotional and physical relationship needed for kids. Schools (primarily) form the intellectual relationships needed.
	
Each division adheres to the community approach, in which larger groups of individuals are always composed of smaller groups.


\subsection{Tiers}
\subsubsection{Tier 1}

3-20 individuals (1 leader) \footnote{While the ideal is 12, I understand that it is hard to have perfect groups)
	Church context: Pastor, church
	School Context: Teacher, class
	Home Context: Parent, family}
\subsubsection{Tier 2}
20-120 individuals\footnote{72 is the desired number for this group, but assuming that Tier 1 is not exactly 12 individuals, 2-6 groups will not equal 72} (2-6 leaders)\footnote{there will always be one leader who is "first among equals" It is important to note that he does not have supreme power, but rather acts as the head for the other leaders. In an abuse of power, other leaders are expected to call him out and rebuke him}

\subsubsection{Tier 3}

120-1440 individuals - 12 groups of a maximum of 120  (6-72 leaders)\footnote{At this point it will be hard to get a unanimous vote, and that is precisely why it must be achieved. Spirit does not give all gifts to all people, and we need to full counsel to best discern the right path. Those who vote a certain way because of sin should be confronted with their wrongdoing.}

\subsubsection{Diagram}
\includegraphics[scale=0.05]{8-social community.png}
This shows that communities are built up from smaller communities. Each small circle is a individual, and is composed of the varying pieces of knowledge that is built up in school.

\section{Technology}
\subsection{}
\subsubsection{Storage and Servers}

\begin{itemize}
\item Quartz glass storage medium\footnote{this ensures that a document can be opened many many years into the future}
\item Must be able to be read easily without a computer - stored in pure text format
\item Local servers that connect with global servers
\item Central Server
\begin{itemize}
	\item All information must be local
	\item All services must be local
	\item provides applications for basic uses\footnote{This does not mean that services outside of a local server cannot be used. This is just a recognition that a teacher should not depend on the Internet for the success of a lesson}
\end{itemize}\end{itemize}

\subsubsection{Formats:}

\emph{All formats used must be compatible by multiple applications}

\begin{itemize}
\item documents (markdown and latex)\footnote{This is not to say any of these technologies are the pinacle of usability, or that in years to come they should be used. Essentially this list represents a single theme: use open software that is free to use, and not locked in to propietary based software. It is expected that better software will emerge under an open liscense, and those new developments should be used}
\item presentations (implress.js, beamer, and markdown)
\item Spreadsheets (Libreoffice Calc)
\item Dvorak as typing system
\item Markdown, Latex, and Pandoc as format for papers
\item Open formats for other
\item OPUS for audio
\item WebM for video
\item WebP for pictures
\item PNG for drawings
\item Git as backend for assignments / lesson plans
\item Open Source Tecnology
\item Linux/Gnu/Ubuntu
\end{itemize}Servers

\subsubsection{Simplicity}

Technology used must be simple enough that anyone with 0-zero knowledge can use it (Policies:0PK Policy)

Some are more capable with technology, and will serve to make the system run simply for others

\subsubsection{Software to use}



\emph{If a open source alternative with similar abilities exist, it will be used}



\begin{itemize}
\item \href{https://stackedit.io/}{StackEdit} for Markdown collaboration
\item Any Latex Editor
\item \href{http://www.ubuntu.com/}{Ubuntu} for servers, tablets, phones, desktops, laptops, and tv's
\item Any code editor
\item \href{http://qt-project.org/}{QT creator} for app design
\item \href{http://www.gimp.org/}{Gimp} for image editing
\item \href{http://www.inkscape.org/en/}{Inkscape} for vector editing
\item \href{./SparkleShare.tex}{SparkleShare} as the client for syncronization
\item Wiki for school organization / curriculum / lesson plans (\href{http://zim-wiki.org/manual/Help/Templates.html}{Zim} does this quite well, but \href{https://www.dokuwiki.org/dokuwiki}{Docuwiki} is another amazing application
\item \href{https://www.gitlab.com/}{GitLab} for collaboration on documents 
\end{itemize}

\end{document}